\chapter{Introduction}
\label{Introduction}

Modern video surveillance systems generate continuous footage at a scale that far exceeds human monitoring capacity. A single security operator watching multiple camera feeds experiences \textit{vigilance decrement} ,  a well-documented decline in attention that sets in after approximately twenty minutes, causing a substantial fraction of visible events to go unnoticed. Automated systems aim to bridge this gap, but they face three persistent challenges. First, visual detectors must track objects across frames to form temporal intervals; without tracking, every frame assigns new identifiers and no interval spans long enough to constitute an event. Second, visual-only systems are blind to events outside the camera's field of view ,  a shout from an adjacent room, a vehicle collision behind a building. Third, naively fusing modalities through temporal alignment reduces recall, because a detection in one modality may fail to find a temporally coincident confirmation in the other.

This thesis addresses these challenges through \textbf{WATCHOUT ISEQL}, a multimodal forensic surveillance framework that combines visual perception from a vision-language model (VLM) with object re-identification and environmental sound analysis from a large audio-language model (LALM). The core of the system is ISEQL (Interval-based Surveillance Event Query Language) \citep{helmer2016highlevel}, which extends relational algebra with temporal operators derived from Allen's interval algebra \citep{allen1983}. The visual stream uses a VLM with a two-phase tracking prompt to maintain consistent object identities across frames, producing relation triples that are collapsed into time intervals. The audio stream uses either Qwen2-Audio-7B-Instruct \citep{chu2024qwen2audio}, a large audio-language model (LALM), or PANNs CNN14 \citep{kong2020panns}, a convolutional neural network pretrained on AudioSet, to produce intervals of detected environmental sounds. Both streams write to a shared SQLite database, and high-level event detection is performed through SQL queries that operate independently on each modality and combine results via set union.

The study is organised as a \textbf{three-condition ablation} experiment. Condition A (visual-only) evaluates VLM detection with object re-identification. Condition B (audio-only) evaluates LALM and CNN-based sound detection across eight window and hop configurations. Condition C (multimodal) takes the union $C = A \cup B$ ,  no temporal cross-modal joins are used, because such joins reduce recall without a compensating benefit. This design allows the contribution of each modality, and of the re-identification mechanism, to be measured in isolation.

This is, to the best of our knowledge, the first application of a Large Audio-Language Model to forensic surveillance event detection. No prior work combining LALMs with interval-based temporal reasoning was found on arXiv or Google Scholar.

\section{Research Aim and Objectives}

The overall aim is to design, implement, and evaluate a multimodal surveillance framework that combines visual and audio perception through a shared interval-based representation, and to quantify the contribution of each modality. The specific objectives are:

\begin{enumerate}
  \item Construct a visual-only pipeline that uses a VLM with object re-identification to produce cross-frame tracking and temporal intervals.
  \item Construct an audio-only pipeline using Qwen2-Audio-7B-Instruct (LALM) and PANNs CNN14 with a closed-vocabulary prompt that constrains the LALM output.
  \item Design a multimodal pipeline that combines visual and audio intervals via union, preserving all detections without temporal cross-modal joins.
  \item Create a labeled evaluation dataset of 30 surveillance scenes annotated with per-frame ground truth for visual relations, audio events, and binary event verdicts.
  \item Evaluate the system across four VLM models, two audio providers across eight window/hop configurations, and a re-identification ablation.
\end{enumerate}

\section{Research Questions}

Three research questions guide the evaluation:

\begin{enumerate}
  \item \textbf{(Object Re-Identification)} How does per-frame object re-identification using class and block-overlap matching affect the number of visual intervals and detected events, compared to the prior approach without re-identification?
  \item \textbf{(Audio Ablation)} How do Qwen2-Audio-7B-Instruct (LALM) and PANNs CNN14 compare across different sliding window and hop configurations on sound-relevant scenes?
  \item \textbf{(Multimodal vs.\ Unimodal)} What is the F1 score of visual-only (A), audio-only (B), and multimodal union (C) across all 30 scenes, and does $C \ge \max(A,B)$ hold on every scene?
\end{enumerate}

\section{Thesis Contributions}

\begin{enumerate}
  \item \textbf{First Large Audio-Language Model for surveillance forensics.} Qwen2-Audio-7B-Instruct is applied to surveillance event detection, achieving perfect precision (1.000) and a micro F1 of 0.788 at the optimal 5.0 s / 2.5 s configuration, substantially outperforming PANNs CNN14 (F1 = 0.429).

  \item \textbf{Object re-identification that enables tracking-dependent events.} A two-phase tracking prompt with class and block-overlap matching is shown to enable the duration-gated events that per-frame ID assignment misses: for the best model (Gemini 3.6 Flash), loitering rises from 0/5 to 5/5, vehicle escape from 0/5 to 1/5, and handoff from 4/5 to 5/5, lifting visual F1 from 0.706 to 0.847.

  \item \textbf{First labeled multimodal forensic surveillance dataset.} A dataset of 30 curated scenes generated with Dreamina Seedance, annotated with per-frame ground truth for visual relations, audio events, and event-level binary verdicts.

  \item \textbf{Union-based multimodal design.} An empirical demonstration that $C = A \cup B$ preserves all detections while temporal cross-modal joins reduce recall.

  \item \textbf{VLM benchmark across four models.} A comparison of Pixtral 12B, Ministral 3-14B, Gemini 2.5 Flash, and Gemini 3.6 Flash on 30 surveillance scenes, establishing Gemini 3.6 Flash as the best performer (F1 = 0.847).

  \item \textbf{Fully auditable forensic pipeline.} Every high-level event is the output of a SQL query that can be read, edited, and re-run ,  no black-box neural network is used in the reasoning path.
\end{enumerate}

\section{Thesis Structure}

The remainder of this thesis is organised as follows. Chapter~\ref{Background} introduces Allen's interval algebra and the ISEQL query language with its general query structure, along with the technical foundations of vision-language models and environmental sound classification. Chapter~\ref{RelatedWork} surveys three strands of prior work: interval-based event detection, VLM-based surveillance, and environmental sound classification. Chapter~\ref{Methodology} presents the three-condition ablation design, the system architecture, the implementation of the visual and audio pipelines, and the formal ISEQL definitions of the six event types. Chapter~\ref{Evaluation} reports the experimental results, structured by research question. Chapter~\ref{Discussion} interprets the findings and discusses threats to validity. Chapter~\ref{ConclusionAndFutureWork} concludes and outlines directions for future work.
